%==============================================================================
% Sjabloon poster bachproef
%==============================================================================
% Gebaseerd op document class `a0poster' door Gerlinde Kettl en Matthias Weiser
% Aangepast voor gebruik aan HOGENT door Jens Buysse en Bert Van Vreckem

\documentclass[a0,portrait,english]{hogent-poster}

\usepackage[table]{xcolor}

% Info over de opleiding
\course{Bachelorproef}
\studyprogramme{toegepaste informatica}
\academicyear{2022-2023}
\institution{Hogeschool Gent, Valentin Vaerwyckweg 1, 9000 Gent}

% Info over de bachelorproef
\title{Server-side Programming Languages for Developing Extensions to IAM Systems}
\subtitle{A Comparative Analysis and Proof of Concept}
\author{Aiko Vandeputte}
\email{aiko.vandeputte@student.hogent.be}
\supervisor{Dhr. J. Willem}
\cosupervisor{Dhr. Y. De Backer (TrustBuilder)}

% Indien ingevuld, wordt deze informatie toegevoegd aan het einde van de
% abstract. Zet in commentaar als je dit niet wilt.
\specialisation{Mobile \& Enterprise developer}
\keywords{Programming languages, Extension, IAM}
\projectrepo{https://github.com/aikovdp/thesis}

\begin{document}

\maketitle

\begin{abstract}
TrustBuilder, an IAM company, wants to rework their extension system, and research which programming language would be best suited to writing these extensions. This thesis performs a comparative analysis for a selection of programming languages, with a proof of concept in Java and TypeScript. A literature review found three main requirements: popularity, memory safety, and static types. Through the proof of concept, both languages were found to be suitable for this use case, but each have their pros and cons. Java appears to have a more mature ecosystem and has a greater focus on correctness. TypeScript on the other hand provides more simplicity, which may be a better fit for a customer-facing extension system. These results can help TrustBuilder make this fundamental decision, and the identified requirements and examples in the proof of concept may guide them in researching the suitability of additional languages.
\end{abstract}

\begin{multicols}{2} % This is how many columns your poster will be broken into, a portrait poster is generally split into 2 columns

\section{Introduction}
TrustBuilder, a European Identity \& Access Management (IAM) company, wants to overhaul its extension system. It is currently part of a monolithic Java application, where JavaScript is used to script what they call Workflows. This system is quite outdated however, and they wish to overhaul it entirely and make it part of their new microservice architecture. To assist them in this task, this thesis researches \textbf{which languages would be the best choice for developing custom extensions for the TrustBuilder IAM platform.} This encompasses the following sub-questions: 
\begin{itemize}
    \item{What are the key indicators these languages should be evaluated against?}
    \item{What would be the pros and cons of specific language choices?}
    \item{What are the security implications of specific language choices?}
\end{itemize}


\section{Methodology}
To provide an answer to these questions, an extensive literature review was first performed to gain insights into IAM standards, comparable platforms, and popular languages and their attributes. A long list of the most popular languages was then evaluated against the requirements gathered through the literature review, to arrive at a short list of 2 languages. These two languages were put to the test in a proof of concept implementing functionality commonly found in TrustBuilder's Workflows. Finally, the findings from this research are summarized in a conclusion, along with guidelines for future research.

\begin{center}
  \captionsetup{type=figure}
  \includegraphics[width=0.4\linewidth]{methodology.png}
  \captionof{figure}{A flowchart representation of the methodology}
\end{center}

\section{Initial requirements analysis and selection}

As customers should be able to write their own extensions, language popularity is an important factor. However, two other notable requirements were found as part of the literature review: memory safety and static typing. Memory safety issues are at the root of many security vulnerabilities, decisively ruling out memory unsafe languages such as C and C++. Static type systems on the other hand help avoid bugs, fix bugs more quickly, and serve as implicit documentation. This certainly makes statically typed languages more appealing candidates.


\begin{center}
  \captionsetup{type=table}
  \begin{tabular}{|l| c c |}
    \hline
    Language & Memory safe & Statically typed \\
    \hline
    JavaScript & Yes & No \\
    Python     & Yes & Certain aspects \\
    Java       & Yes & Yes \\
    TypeScript & Yes & Yes \\
    C++        & \textcolor{red}{\textbf{No}}  & Yes \\
    C\#        & Yes & Yes \\
    C          & \textcolor{red}{\textbf{No}}  & Yes \\
    PHP        & Yes & Certain aspects \\
    Go         & Yes & Yes \\
    Kotlin     & Yes & Yes \\
    Rust       & Yes & Yes \\
    \hline
\end{tabular}
  \captionof{table}{Requirements summary for popular server-side programming languages, ordered by descending popularity}
\end{center}

Checking the long list of the most popular programming languages against memory safety and type safety requirements, Java and TypeScript emerge as the most promising options to explore further in a proof of concept. These are also the two languages used by TrustBuilder's own development teams, making them even more valuable options.

\section{Conclusions}


As was expected, both TypeScript and Java are excellent choices, but they both have their pros and cons. Java proved to be more reliable, with a more mature ecosystem and a bigger focus on safety. TypeScript offered more simplicity however, which may be better suited to a customer-facing extension system, though more care would have to be taken in choosing the tech stack.

\section{Future research}

Both TypeScript and Java are fine choices for TrustBuilder’s case, but others may be excellent as well. This thesis focused on these two because of their popularity and the existing internal experience with them at TrustBuilder. Following the requirements mentioned earlier and a similar proof of concept, TrustBuilder could also verify how other languages such as Kotlin, Rust or Go might function in their next generation of workflows.

\end{multicols}
\end{document}